% ===================================================================
%  TABULKA č. 14 — Ulice. — Obchodníci.
%  Traduction 2026 d'apres texto/io, controlee sur texto/fr, texto/ru
%  et texto/uk. Consigne : voir texto/cs/10-tabulka-01.tex.
%
%  LE TABLEAU DES METIERS EST LE PLUS LONG DE LA COLONNE, et il donne
%  ici la meme mesure qu'en irlandais et en galicien : LE MOT A L'AGE
%  DE SON OBJET. Ce qui se fait de ses mains depuis toujours garde
%  son mot tcheque — « kovář » le forgeron, « holič » le coiffeur,
%  « krejčí » le tailleur, « zlatník » l'orfevre, « klempíř » le
%  ferblantier, « čalouník » le tapissier, « kominík » le ramoneur,
%  « podomní obchodník » le colporteur — et ce qui s'est invente au
%  XIXe siecle arrive avec son nom : « fotograf », « litograf »,
%  « stenograf », « fonograf », « optik », « cyklista ».
%
%  MAIS LE TCHEQUE FAIT ICI CE QU'AUCUNE AUTRE COLONNE N'A FAIT : il
%  refait le suffixe lui-meme. La ou le galicien prenait le grec en
%  bloc — fotógrafo, litógrafo, mecanógrafo — le tcheque garde
%  « fotograf » et « litograf » mais dit « písař na stroji » pour
%  la dactylographe, « knihkupec » pour le libraire, « knihař » pour
%  le relieur, « tiskař » pour l'imprimeur. Le grec entre quand
%  l'objet est neuf ; il ne franchit pas la porte d'un metier qui
%  existait deja sous un autre nom.
% ===================================================================

%%K t14-apar-1 apar
\VUsaut{4.0mm}
\VUtitre{15.0pt}{60}{TABULKA č. 14}
\VUsaut{3.0mm}

%%K t14-tit-1 sub
\VUpk{11.4pt}{40}{Ulice. --- Obchodníci.}
\VUsaut{3.0mm}

%%K t14-01-1 p
1. --- Můj přítel Jan, který bydlí na venkově, přijel strávit tři dny v mé rodině.
\VUgras{Až} dosud neměl příležitost vidět velké město; proto trávíme všechny své dny
procházkami a já mu vysvětluji, co nezná.

%%K t14-02-1 p
2. --- Nejprve jsme šli navštívit \VUgras{hvězdárnu} \textsuperscript{(1)}, která ho
velmi zajímala. Když jsme se vraceli \VUgras{ulicí} \textsuperscript{(3)}, kde jsou
\VUgras{turecké lázně} \textsuperscript{(2)}, potkali jsme \VUgras{pohřeb}
\textsuperscript{(4)}. Před \VUgras{pohřebním průvodem} kráčelo \VUgras{duchovenstvo,}
které předcházelo \VUgras{pohřební vůz,} v němž byla položena \VUgras{rakev.} Mnoho
přátel doprovázelo \VUgras{truchlící} rodinu.

%%K t14-02-2 p
Když průvod prošel, přešli jsme ulici před velkým \VUgras{pivovarským hostincem}
\textsuperscript{(5)}, v němž mnoho \VUgras{hostů} pilo \VUgras{pivo} na terase,
chráněno prosklenou \VUgras{markýzou} \textsuperscript{(6)}.

%%K t14-03-1 p
3. --- Jan byl velmi udiven \VUgras{štítem} umístěným mezi obchodem
\VUgras{lihovarníka} \textsuperscript{(7)} a obchodem \VUgras{prodavače hudebnin}
\textsuperscript{(10)}. Zeptal se mě na jeho význam; odpověděl jsem, že označuje
anglický \VUgras{konzulát} \textsuperscript{(8)}, a upozornil jsem ho na
\VUgras{konzula} \textsuperscript{(9)}, který je na balkoně.

%%K t14-tit-2 sub
\VUpk{10.2pt}{40}{Prohlídka obchodů.}
\VUsaut{3.0mm}

%%K t14-04-1 p
4. --- Ovšem jsme se zastavili před výkladem \VUgras{hudebních nástrojů,}
\VUgras{harmonií,} \VUgras{varhan} a \VUgras{klavírů}; dívali jsme se na ilustrované
obálky \VUgras{partitur} a \VUgras{hudebnin} pro klavír a obdivovali jsme
\VUgras{lyru} ze zlaceného dřeva, které se užívá jako štítu.

%%K t14-05-1 p
5. --- Oblaka prachu, která zvířili \VUgras{zametači} \textsuperscript{(11)} a
\VUgras{zametací vůz} \textsuperscript{(12)}, nás přinutila projít rychle před krásným
\VUgras{nábytkem} \VUgras{obchodníka s nábytkem} \textsuperscript{(13)} a před
výkladem \VUgras{kamen} a \VUgras{lamp} \VUgras{klempíře} \textsuperscript{(14)}.

%%K t14-06-1 p
6. --- Ostatně na tom místě jsme byli nuceni opustit chodník, neboť byl zatarasen
balíky, které se braly z \VUgras{dopravního vozu} \textsuperscript{(16)}, aby se
uložily do \VUgras{obchodu} \textsuperscript{(15)}, který si právě pronajali.

%%K t14-06-2 p
To nám zabránilo vidět krásné vozy \VUgras{kočárníka} \textsuperscript{(17)}, ale
nezabránilo nám zastavit se a podívat se na \VUgras{balírníka}
\textsuperscript{(19)}, jak dělá bednu pro svého souseda, \VUgras{obchodníka s koly a
automobily} \textsuperscript{(18)}. Potom, protože už bylo poněkud pozdě, vrátili jsme
se domů.

%%K t14-tit-3 sub
\VUpk{10.2pt}{40}{Procházka městem.}
\VUsaut{3.0mm}

%%K t14-07-1 p
7. --- Druhého dne navrhla má matka Janovi, aby se dal vyfotografovat, a potom mi
uložila, abych ho doprovodil k \VUgras{fotografovi} \textsuperscript{(20)}. Po obědě
jsme šli do \VUgras{ateliéru} umělcova, kde jsme byli nuceni dosti dlouho čekat.
Abychom se zabavili, dívali jsme se na \VUgras{podobizny} \textsuperscript{(22)}
zavěšené na stěně.

%%K t14-07-2 p
Když přišla naše řada, práce netrvala dlouho, a jakmile můj přítel přestal pózovat před
\VUgras{fotografickým přístrojem} \textsuperscript{(21)}, sešli jsme dolů.

%%K t14-08-1 p
8. --- V prvním patře jsme viděli dveřmi \VUgras{holiče} \textsuperscript{(23)}, které
jsou otevřené, jeho pomocníka, jak stříhá vlasy zákazníkovi.

%%K t14-08-2 p
Když jsme přišli na ulici, prošli jsme před \VUgras{trafikou} \textsuperscript{(24)}.
Jana tak pobavila červená \VUgras{lucerna} \textsuperscript{(25)} a \VUgras{tlustá
dýmka}, které se užívá jako \VUgras{štítu} \textsuperscript{(26)}, že vrazil do dvou
starých pánů, kteří rozmlouvali \VUgras{šňupajíce tabák} \textsuperscript{(27)}.

%%K t14-tit-4 sub
\VUpk{10.2pt}{40}{Cukrárna.}
\VUsaut{3.0mm}

%%K t14-09-1 p
9. --- \VUgras{Bankovky} a \VUgras{mince} ze všech zemí vystavené ve \VUgras{výkladu}
\VUgras{směnárníka} \textsuperscript{(29)} na chvíli poutaly naši pozornost, ale
protože byla hodina svačiny, vstoupili jsme k \VUgras{cukráři}
\textsuperscript{(31)}, aniž jsme se déle zdrželi.

%%K t14-10-1 p
10. --- Když jsme viděli všechny \VUgras{pochoutky,} které obchod obsahoval,
\VUgras{koláče, perník, sušenky} a všeliké \VUgras{bonbony,} nebylo pro nás snadné
vybrat si. Nakonec si Jan vybral \VUgras{krémové koláčky} a já jsem vzal jen malý
\VUgras{třešňový dortík,} neboť sladkosti mi \VUgras{kazí zuby,} a nechci vůbec chodit
příliš často k \VUgras{zubaři} \textsuperscript{(30)}.

%%K t14-tit-5 sub
\VUpk{10.2pt}{40}{Psaní na stroji.}
\VUsaut{3.0mm}

%%K t14-11-1 p
11. --- Abych svému příteli ukázal \VUgras{psací stroj}, o němž slyšel mluvit, ale
který neznal, vstoupili jsme do \VUgras{kanceláře} \textsuperscript{(32)} svého
\VUgras{bratrance} \textsuperscript{(34)}. Nalezli jsme ho, jak diktuje své dopisy svému
\VUgras{tajemníkovi} \textsuperscript{(35)}, který je \VUgras{stenograf}. Sotva byl
dopis dokončen, \VUgras{písař na stroji} \textsuperscript{(33)} jej opsal na svém
stroji. Protože jsme nechtěli přerušovat zaměstnání svého příbuzného, zůstali jsme u
něho jen několik okamžiků.

%%K t14-12-1 p
12. --- Pod kanceláří mého bratrance bydlí velký \VUgras{hodinář-klenotník}
\textsuperscript{(37)}, který je nadto zlatník. Jeho skvělý obchod obsahuje mnoho
\VUgras{hodin, kyvadlových hodin, kapesních hodinek,} \VUgras{řetízků} a drahocenných
\VUgras{klenotů,} například \VUgras{perlových náhrdelníků,} zlatých \VUgras{náramků,}
\VUgras{náušnic} a \VUgras{prstenů} zdobených \VUgras{drahokamy} a \VUgras{diamanty}.
Jeden z výkladů je zvlášť vyhrazen \VUgras{zlatnickému zboží} a obsahuje významnou
sbírku stříbrných \VUgras{příborů}.

%%K t14-13-1 p
13. --- Když jsme zahýbali za roh ulice, všiml jsem si, že vyměnili \VUgras{tabulku}
\textsuperscript{(36)} umístěnou blízko \VUgras{čísla} domu. Chtěl jsem svou poznámku
říci nahlas, když zazněly veselé zvuky \VUgras{vojenské} hudby. Rozběhli jsme se
vstříc pluku, aniž jsme uvážili, že \VUgras{ten} má projít před obchody
\VUgras{kloboučníka} \textsuperscript{(39)} a \VUgras{rukavičkáře}
\textsuperscript{(40)}, a že bychom byli právě tak dobře umístěni, abychom viděli jeho
přehlídku, kdybychom zůstali před výkladem \VUgras{optika} \textsuperscript{(38)},
zcela zásobeným \VUgras{skřipci, lorňony} a \VUgras{kukátky}.

%%K t14-tit-6 sub
\VUpk{10.2pt}{40}{Přehlídkový pochod.}
\VUsaut{3.0mm}

%%K t14-14-1 p
14. --- Před \VUgras{plukem} \textsuperscript{(41)} kráčel \VUgras{tambormajor}
\textsuperscript{(42)}, následovaný \VUgras{bubeníky} \textsuperscript{(43)},
\VUgras{trubači} \textsuperscript{(44)} a celou \VUgras{kapelou}. \VUgras{Generál}
\textsuperscript{(45)}, který byl na náměstí se svým pobočníkem, se zastavil blízko
\VUgras{vodotrysku} \textsuperscript{(49)} \VUgras{kašny} \textsuperscript{(48)}, aby
byl přítomen \VUgras{přehlídce.}

%%K t14-14-2 p
\VUgras{Štábní důstojníci} \textsuperscript{(46)}, \VUgras{plukovník} a
\VUgras{podplukovník} ho pozdravili kordem. \VUgras{Majoři} byli před svými
\VUgras{prapory} \textsuperscript{(47)} a každý \VUgras{kapitán} velel své
\VUgras{rotě,} zatímco \VUgras{poručíci,} \VUgras{podporučíci} a \VUgras{poddůstojníci}
zaujímali své obvyklé místo vedle svých mužů.

%%K t14-15-1 p
15. --- Mnoho chodců se shluklo na \VUgras{křižovatce} \textsuperscript{(50)};
obchodníci, \VUgras{deštníkář} \textsuperscript{(51)}, \VUgras{cínař}
\textsuperscript{(53)} a \VUgras{puškař} \textsuperscript{(54)} vyšli, aby viděli
\VUgras{vojáky.} Všechna vozidla, \VUgras{fiakry} \textsuperscript{(52)},
\VUgras{stěhovací vozy} \textsuperscript{(57)} i \VUgras{kropicí vůz}
\textsuperscript{(56)}, se postavila podél chodníku.

%%K t14-tit-7 sub
\VUpk{10.2pt}{40}{Výjevy na ulici.}
\VUsaut{3.0mm}

%%K t14-15-2 p
\VUgras{Obchodní cestující} \textsuperscript{(55)} nesoucí v ruce své \VUgras{krabice
se vzorky} rychle přešel ulici a osoby sedící na \VUgras{imperiálu}
\textsuperscript{(59)} \VUgras{omnibusu} \textsuperscript{(58)} se obrátily, aby si
užily podívané. Ubohý \VUgras{potulný zpěvák} \textsuperscript{(60)} se svým
\VUgras{flašinetem} \textsuperscript{(61)} musel přerušit svou \VUgras{písničku,}
neboť hluk bubnů přehlušil jeho hlas.

%%K t14-16-1 p
16. --- Když pluk došel před \VUgras{obchod s látkami} \textsuperscript{(64)},
\VUgras{švadleny} \textsuperscript{(63)} a \VUgras{krejčí} \textsuperscript{(62)}
opustili své \VUgras{šicí stroje} a svou práci, aby se podívali oknem.
\VUgras{Obchodník} \textsuperscript{(65)}, který právě položil nové \VUgras{obleky}
\textsuperscript{(66)} na svůj \VUgras{výklad,} musel dát dovnitř \VUgras{sukna}
\textsuperscript{(67)} a \VUgras{plátna} vyložená na chodníku, blízko
\VUgras{kanálového otvoru} \textsuperscript{(68)}, ze strachu, že je zástup převrhne;
dokonce i starý \VUgras{podomní obchodník} \textsuperscript{(69)}, s nímž domovnice
smlouvala o punčochy, byl nucen vejít do průjezdu, aby dokončil svůj prodej.

%%K t14-16-2 p
Doprovodili jsme pluk až ke kasárnám \VUgras{a,} velmi spokojeni se svým dnem, vrátili
jsme se v hodinu večeře.

%%K t14-tit-8 sub
\VUpk{10.2pt}{40}{Různé obchody a povolání.}
\VUsaut{3.0mm}

%%K t14-17-1 p
17. --- Následujícího dne nám otec navrhl navštívit tiskárnu tamních novin. S nadšením
jsme přijali.

%%K t14-17-2 p
\VUgras{Tiskař} je také \VUgras{knihkupec} \textsuperscript{(70)} \VUgras{(= prodavač
knih),} \VUgras{nakladatel,} \VUgras{papírník} a \VUgras{knihař}. V prvním patře zřídil
\VUgras{redakci} \textsuperscript{(71)} novin a také \VUgras{rytecký ateliér,} v němž
pracují \VUgras{litografové} \textsuperscript{(72)} a \VUgras{rytci do mědi}
\textsuperscript{(73)}, konečně \VUgras{sazárnu,} v níž jsou \VUgras{tiskařští
dělníci} \textsuperscript{(74)}. Vlastní \VUgras{tiskárna} \textsuperscript{(75)},
\VUgras{tiskařské lisy} a různé stroje jsou v přízemí.

%%K t14-tit-9 sub
\VUpk{10.2pt}{40}{Jan klade velmi mnoho otázek.}
\VUsaut{3.0mm}

%%K t14-18-1 p
18. --- Když jsme vycházeli z tiskárny, Jan se velmi udiven zastavil před malou
prosklenou skříňkou umístěnou na sloupu, proti \VUgras{obchodu s galanterií}
\textsuperscript{(77)}, a ptal se, k čemu slouží. Můj otec mu vysvětlil, že je to
\VUgras{požární hlásič} \textsuperscript{(76)}. Ostatně všechno ve městě bylo pro mého
přítele předmětem údivu, od prosté \VUgras{kamenné kašny} \textsuperscript{(78)} a
lehké \VUgras{nákladní tříkolky} \textsuperscript{(80)} řízené \VUgras{poslíčkem}
\textsuperscript{(79)} v livreji až po \VUgras{poštovní vůz} \textsuperscript{(81)}.

%%K t14-19-1 p
19. --- Zatímco nás otec na chvíli opustil, aby promluvil s \VUgras{výběrčím daní}
\textsuperscript{(83)}, který se zastavil, aby rozmlouval s \VUgras{advokátem}
\textsuperscript{(82)} a s \VUgras{notářem} \textsuperscript{(84)}, bavili jsme se
sledováním ruchu na ulici. Míjely nás nejrozmanitější osoby: \VUgras{cukrářský
pomocník} \textsuperscript{(85)}, nesoucí v koši skvělý \VUgras{paštikový koláč,} šel
za \VUgras{obchodníkem s oděvy} \textsuperscript{(86)}; \VUgras{lampář}
\textsuperscript{(87)} rozmlouval s \VUgras{plynařem} \textsuperscript{(88)};
\VUgras{řeholnice} \textsuperscript{(89)}, která vedla za ruku sirotu, se vracela do
svého kláštera.

%%K t14-tit-10 sub
\VUpk{10.2pt}{40}{Při našem návratu domů.}
\VUsaut{3.0mm}

%%K t14-20-1 p
20. --- V okamžiku, kdy nás otec dohonil, Jan se hlasitě rozesmál, když uviděl špinavou
tvář a podivné oblečení dvou \VUgras{kominíků} \textsuperscript{(91)} zcela černých od
sazí.

%%K t14-21-1 p
21. --- Chtěl jsem koupit své matce rostlinu od \VUgras{květinářky}
\textsuperscript{(92)}, ale otec mě upozornil, že by byla velmi těžká na nesení a že by
bylo lepší pořídit \VUgras{pomeranče}. Přiblížili jsme se k \VUgras{prodavačce}
\textsuperscript{(94)}, a když \VUgras{vychovatelka} \textsuperscript{(93)}, která
vybírala pro svou svěřenku, dokončila svůj nákup, dali jsme se obsloužit.

%%K t14-22-1 p
22. --- Když jsme se vraceli domů, potkali jsme několik známých: starou přítelkyni mé
rodiny, která se opírala paží o svou \VUgras{společnici} \textsuperscript{(95)},
\VUgras{inženýra} \textsuperscript{(96)} mostů a silnic a jednu ze svých sestřenic s
její \VUgras{chůvou} \textsuperscript{(98)}.

%%K t14-22-2 p
Potom jsme potkali \VUgras{modistku} \textsuperscript{(97)} mé matky a dva
\VUgras{mnichy} \textsuperscript{(99)}, kteří byli ve městě cizí a přistoupili k nám,
aby se zeptali mého otce na cestu k nádraží.
